\documentclass[10pt,a4paper]{amsart}
\usepackage[margin=19mm]{geometry}
\usepackage{amsmath,amssymb,mathtools}
\usepackage[scr=rsfs,cal=euler]{mathalfa}
\usepackage{xfrac}
\usepackage[unicode,hidelinks,bookmarks=false]{hyperref}
\newcommand{\ie}{i.e.\ }
\newcommand{\cf}{cf.\ }
\newcommand{\resp}{respectively\ }
\newcommand{\Subsection}{\subsection}
\providecommand{\nobreakdash}{\nobreak\hbox{-}}
\setlength{\emergencystretch}{2em}
\multlinegap0pt
\usepackage[shortlabels]{enumitem}
\usepackage{amssymb}
\usepackage{mathtools}
\newtheorem{theorem}{Theorem}
\newcommand{\BM}{\mathrm{BM}}
\newcommand{\dist}{d_{\BM}}
\renewcommand{\le}{\leqslant}
\title{$(\lambda^+)$-injective Banach spaces}
\author{Tomasz Kania, Grzegorz Lewicki}
\date{Source: arXiv:2603.09710v1 (CC BY 4.0)}
\begin{document}
\maketitle

\noindent\emph{Source and licence.} $(\lambda^+)$-injective Banach spaces. Version arXiv:2603.09710v1 (CC BY 4.0), distributed by arXiv under the Creative Commons Attribution 4.0 International licence, http://creativecommons.org/licenses/by/4.0/, as recorded in the arXiv licence metadata for this exact version and shown on the abstract page https://arxiv.org/abs/2603.09710v1. Full licence notice: \texttt{licenses/arxiv-2603.09710-CC-BY-4.0.txt}.\par\smallskip

\noindent\textit{Editorial scope.} Counted unit: the article's theorem thm:optimised (source0.tex lines 521--530). The article's complete original proof of it is delivered with it (source0.tex lines 532--674, one \texttt{proof} environment of 143 lines). No mathematics is written, repaired or re-derived: the statement and the proof are the article's own lines, in the article's own notation, and no retained line is substituted or re-flowed.  No further source-local context is needed: every cross-reference of every retained line resolves inside the two delivered spans.  Nothing was omitted from any retained span, so the package carries no omission record.  No retained line contains a citation, so the wrapper carries no bibliography and no bibliography entry.  No retained line contains a figure environment, an image-inclusion command or drawing code, so no image is delivered and the package declares no asset dependency.  Editorial formatting only, in addition: an AMS \texttt{amsart} preamble that restates the article's own declarations and macros used by the retained lines, namely the environments \texttt{theorem} on separate counters, together with the standard packages those lines need.  The retained environments print in this document as Theorem 1 in order of appearance, whereas the article prints the same items with the numbering of its own full document.  The complete native source of the article as distributed in the arXiv e-print for this exact version is bundled unchanged as \texttt{sources/196083/source0.tex}.
\begin{theorem}\label{thm:optimised}
Let $X$ and $Y$ be Banach spaces such that
\begin{enumerate}[label=\textup{(\roman*)}]
\item there exist $1$-complemented subspaces $X'\subseteq X$ and
$Y'\subseteq Y$ with $X$ isometric to $Y'$ and $Y$ isometric to $X'$;
\item $X$ is isometric to $X\oplus_\infty X$ and $Y$ is isometric
to $Y\oplus_\infty Y$.
\end{enumerate}
Then $\dist(X,Y)\le 9+6\sqrt{3}$.
\end{theorem}

\begin{proof}
Fix norm-one projections $P\colon X\to X'$ and $R\colon Y\to Y'$, and
set $E:=\ker P$ and $F:=\ker R$.  By~(i), fix surjective isometries
$\theta\colon X\to Y'$ and $\eta\colon Y\to X'$.  By~(ii), fix
surjective isometries
$\varphi=(\varphi_1,\varphi_2)\colon X\to X\oplus_\infty X$ and
$\psi=(\psi_1,\psi_2)\colon Y\to Y\oplus_\infty Y$.

Let $a>0$ and put
\[
\mu:=\frac{1}{a},\qquad
\nu:=\frac{\sqrt{2a+1}}{a},\qquad
b:=\frac{a}{\sqrt{2a+1}}=\frac{1}{\nu}.
\]
Define three operators:
\begin{align*}
T_a\colon X&\to Y\oplus_\infty Y\oplus_\infty E,\\
T_a x&:=\bigl(\nu\,\psi_1\eta^{-1}Px,\;
\mu\,\psi_2\eta^{-1}Px,\; x-Px\bigr);\\[2mm]
S\colon Y\oplus_\infty Y\oplus_\infty E
&\to X\oplus_\infty X\oplus_\infty F,\\
S(y_1,y_2,e)&:=\bigl(\theta^{-1}Ry_1,\;
\eta y_2+e,\; y_1-Ry_1\bigr);\\[2mm]
U_a\colon X\oplus_\infty X\oplus_\infty F&\to Y,\\
U_a(x_1,x_2,f)&:=\theta\,\varphi^{-1}(ax_1,\,bx_2)+f.
\end{align*}
Each is bijective, with inverses:
\begin{align*}
T_a^{-1}(y_1,y_2,e)
&=\eta\,\psi^{-1}\!\bigl(\tfrac{1}{\nu}y_1,\,ay_2\bigr)+e,\\
S^{-1}(x_1,x_2,f)
&=\bigl(\theta x_1+f,\;\eta^{-1}Px_2,\;x_2-Px_2\bigr),\\
U_a^{-1}y
&=\bigl(\tfrac{1}{a}\,\varphi_1\theta^{-1}Ry,\;
\nu\,\varphi_2\theta^{-1}Ry,\;y-Ry\bigr).
\end{align*}
Set $W_a:=U_aST_a\colon X\to Y$.

\smallskip
\noindent\emph{Estimating $\|W_a\|$.}
Let $\|x\|\le 1$.  Since $\|Px\|\le 1$ and $\|x-Px\|\le 2$,
and $\|\psi_1\eta^{-1}Px\|\le 1$, while
$\|(I_Y-R)\psi_1\eta^{-1}Px\|\le 2$, a direct expansion yields
$\|W_ax\|\le\max\{a\nu,\,b(\mu+2)\}+2\nu$.
By the choice of parameters, $a\nu=\sqrt{2a+1}$ and
$b(\mu+2)=\sqrt{2a+1}$, so
$\|W_a\|\le 2\nu+\sqrt{2a+1}$.

\smallskip
\noindent\emph{Estimating $\|W_a\|$.}
Let $\|x\|\le 1$.  Writing out the composition, we obtain
\[
W_ax
=
\theta\varphi^{-1}\!\Bigl(
a\nu\,\theta^{-1}R\psi_1\eta^{-1}Px,\;
b\bigl(\mu\,\eta\psi_2\eta^{-1}Px+x-Px\bigr)
\Bigr)
+
\nu(I_Y-R)\psi_1\eta^{-1}Px.
\]
Hence
\begin{align*}
\|W_ax\|
&\le
\max\Bigl\{
a\nu\|R\psi_1\eta^{-1}Px\|,
\,b\bigl(\mu\|\psi_2\eta^{-1}Px\|+\|x-Px\|\bigr)
\Bigr\}
+\nu\|(I_Y-R)\psi_1\eta^{-1}Px\| \\
&\le \max\{a\nu,\,b(\mu+2)\}+2\nu,
\end{align*}
because $\|Px\|\le 1$, $\|x-Px\|\le 2$, $\|R\|=1$, and
$\|(I_Y-R)\|\le 1+\|R\|=2$.
By the choice of parameters,
\[
a\nu=\sqrt{2a+1}
\qquad\text{and}\qquad
b(\mu+2)=\sqrt{2a+1},
\]
so
\[
\|W_a\|\le 2\nu+\sqrt{2a+1}.
\]

\smallskip
\noindent\emph{Estimating $\|W_a^{-1}\|$.}
Let $\|y\|\le 1$.  Since $W_a^{-1}=T_a^{-1}S^{-1}U_a^{-1}$, we have
\[
U_a^{-1}y
=
\bigl(
\mu\,\varphi_1\theta^{-1}Ry,\;
\nu\,\varphi_2\theta^{-1}Ry,\;
y-Ry
\bigr),
\]
and therefore
\[
S^{-1}U_a^{-1}y
=
\bigl(
\mu\,\theta\varphi_1\theta^{-1}Ry+y-Ry,\;
\nu\,\eta^{-1}P\varphi_2\theta^{-1}Ry,\;
\nu(I_X-P)\varphi_2\theta^{-1}Ry
\bigr).
\]
Applying $T_a^{-1}$ gives
\[
W_a^{-1}y
=
\eta\psi^{-1}\!\Bigl(
\frac{1}{\nu}\bigl(\mu\,\theta\varphi_1\theta^{-1}Ry+y-Ry\bigr),\;
a\nu\,\eta^{-1}P\varphi_2\theta^{-1}Ry
\Bigr)
+
\nu(I_X-P)\varphi_2\theta^{-1}Ry.
\]
Consequently,
\begin{align*}
\|W_a^{-1}y\|
&\le
\max\Bigl\{
\frac{1}{\nu}(\mu+2),\;
a\nu
\Bigr\}
+2\nu \\
&=
\max\{b(\mu+2),\,a\nu\}+2\nu
=
2\nu+\sqrt{2a+1}.
\end{align*}
Thus $\|W_a^{-1}\|\le 2\nu+\sqrt{2a+1}$.

\smallskip
Setting $K(a):=2\nu+\sqrt{2a+1}=\frac{(a+2)}{a}\sqrt{2a+1}$,
we obtain $\dist(X,Y)\le K(a)^2=g(a)$ where
\[
g(a):=2a+9+\frac{12}{a}+\frac{4}{a^2}.
\]
Solving $g'(a)=0$ yields $a^3-6a-4=0$, whose unique positive root is
$a=1+\sqrt{3}$.  Substituting gives $g(1+\sqrt{3})=9+6\sqrt{3}$.
\end{proof}

\end{document}
